\section{How the physics works}\label{sec:physics}

\begin{inshort}
Pumping creates a cone of low water pressure around a well field. Clay inside that cone compacts, and
the ground above sinks in the same cone-shaped pattern. Our inversion reads that pattern backwards.
\end{inshort}

\begin{figure}[H]
\centering
\begin{tikzpicture}[node distance=5mm]
\tikzset{ch/.style={box,text width=24mm,minimum height=12mm}}
\node[ch,boxw] (a) {Pumping\\{\scriptsize volume per year}};
\node[ch,right=of a] (b) {Water pressure\\drops in a cone};
\node[ch,right=of b] (c) {Clay layers\\compact};
\node[ch,right=of c] (d) {Ground\\sinks};
\node[ch,boxa,right=of d] (e) {Radar measures\\the sinking};
\draw[arr] (a)--(b); \draw[arr] (b)--(c); \draw[arr] (c)--(d); \draw[arr] (d)--(e);
\draw[arr,draw=accent,line width=0.9pt] (e.south) -- ++(0,-5mm) -| (a.south);
\node[lab,text=accent] at ($(c.south)+(0,-9.5mm)$) {our inversion runs backwards along this chain};
\node[lab] at ($(c.north)+(0,3mm)$) {physics runs forwards};
\end{tikzpicture}
\caption{The causal chain. Each arrow is one piece of physics described below.}
\end{figure}

\subsection{From water pressure to compaction}

When the water level (head) $h$ in an aquifer system falls, a layer of thickness $b_0$ compacts by
\begin{equation}
\Delta b \;=\; S_k\,\Delta h, \qquad
S_k = \begin{cases} S_{ke} & \text{if } h \text{ is above its lowest past level (elastic, recoverable)}\\
                    S_{kv} & \text{if } h \text{ falls below it (inelastic, permanent)} \end{cases}
\label{eq:compaction}
\end{equation}
where $S_k$ is the skeletal storage coefficient. $S_{kv}$ is typically one to two orders of magnitude
larger than $S_{ke}$ \cite{galloway2011}. This is why over-pumping is so damaging: once head falls
below its historical minimum, compaction jumps and most of it never recovers. In Semarang, the Damar
aquifer's head has already dropped below that threshold \cite{lo2022}.

Thick clay drains slowly. A clay layer of thickness $b'$, vertical conductivity $K'_v$ and specific
storage $S'_s$ needs a time of the order of
\begin{equation}
\tau \;=\; \frac{(b'/2)^2\,S'_s}{K'_v}
\end{equation}
to adjust \cite{galloway2011}. For metres of clay this is years to decades. Two consequences follow.
The ground keeps sinking long after pumping stabilises, and it sinks in the place where the pressure
cone sits. The second point is what makes attribution possible.

\subsection{Why an elastic point source is not enough}\label{sec:identity}

The usual model for inverting ground motion is a point source of volume change $\Delta V$ at depth
$d$ in an elastic half-space (the Mogi model \cite{mogi1958,geertsma1973}). Its vertical surface
motion at horizontal distance $r$ is
\begin{equation}
u_z(r) \;=\; \frac{(1-\nu)\,\Delta V}{\pi}\;\frac{d}{(r^2+d^2)^{3/2}},
\label{eq:mogi}
\end{equation}
with Poisson's ratio $\nu$. Azeriansyah et al.\ used this model for Semarang with $d = 120$\,m
\cite{azeriansyah2025}. Two short calculations show what it can and cannot tell us.

\textbf{Total volume.} Integrating \eqref{eq:mogi} over the surface gives
$\int u_z\,\mathrm{d}A = 2(1-\nu)\,\Delta V$. The model therefore returns a volume that is simply the
integrated subsidence times a constant.

\textbf{Locality.} The share of a source's surface signal that falls within radius $R$ of it is
\begin{equation}
f(R) \;=\; 1-\frac{d}{\sqrt{R^2+d^2}} .
\end{equation}
For $d = 120$\,m this is 57\,\% within 250\,m and 77\,\% within 500\,m. A source at aquifer depth
moves the ground almost only directly above itself.

\begin{plain}
At a depth of about 100\,m, an elastic point source is nearly a mirror: the inverted ``volume map''
is the subsidence map again, rescaled. It cannot say whether a 3\,km wide bowl comes from one big
well field in the middle or from thirty small wells spread across it. Something else must carry that
information.
\end{plain}

\subsection{The missing piece: the shape of a pumping cone}

That ``something else'' is hydraulics. Around a well field that pumps at rate $Q$ from a confined
aquifer with leaky clay above it, the long-term drawdown is (Hantush and Jacob \cite{hantush1955})
\begin{equation}
s(r) \;=\; \frac{Q}{2\pi T}\;K_0\!\left(\frac{r}{B}\right), \qquad B=\sqrt{\frac{T\,b'}{K'_v}} ,
\label{eq:hantush}
\end{equation}
where $T$ is the aquifer's transmissivity, $K_0$ the modified Bessel function, and $B$ the
\emph{leakage factor}: the distance over which the cone fades. Compaction rate follows drawdown
(Eq.~\ref{eq:compaction}), so each source $j$ at position $\mathbf{x}_j$ adds a cone-shaped piece to
the vertical velocity field:
\begin{equation}
v(\mathbf{x}) \;=\; \sum_{j} q_j\,k_B\big(\lVert\mathbf{x}-\mathbf{x}_j\rVert\big) \;+\; g(\mathbf{x}) \;+\; \varepsilon(\mathbf{x}),
\qquad k_B(r)=\frac{K_0(r/B)}{2\pi B^2}.
\label{eq:model}
\end{equation}
The kernel $k_B$ integrates to one, so $q_j$ is directly the \textbf{compaction volume per year}
caused by source $j$ (m$^3$/yr). $g$ is the non-pumping background, and $\varepsilon$ is measurement
noise. Equation~\eqref{eq:model} is the forward model of this project.

\begin{caution}
The cone shape is a modelling choice. Equation~\eqref{eq:hantush} assumes a uniform aquifer and a
steady cone. Real aquifers are patchy, and $B$ is not known in advance. Section~\ref{sec:exp} measures
how wrong the answers get when the cone shape or width is wrong, and Section~\ref{sec:feas} says how
to calibrate $B$.
\end{caution}

\subsection{From compaction volume to extraction}\label{sec:massbalance}

Pumped water comes from three places: water squeezed out of compacting clay, water released
elastically, and water that leaks in from above or from the sides. Only the first leaves a permanent
mark on the ground. So for every source
\begin{equation}
Q_j \;=\; \frac{q_j}{\varphi_j}, \qquad 0<\varphi_j\le 1 ,
\label{eq:phi}
\end{equation}
where $\varphi_j$ is the share of pumped water that comes from compaction. Because
$\varphi_j \le 1$, \textbf{the compaction volume $q_j$ is a lower bound on extraction}. Vasco et
al.\ reached the same conclusion for California \cite{vasco2019}. With metered wells or head records
(\D{6}), $\varphi$ can be calibrated and the bound becomes an estimate.

\textbf{Order-of-magnitude check.} Suppose 75\,km$^2$ sinks at an average of 5\,cm/yr. (The area is
20\,\% of the city \cite{aditiya2023}; the average rate is our assumption.) That is about
3.7 million m$^3$ of lost storage per year, against more than 30 million m$^3$ pumped
\cite{putranto2016}. This gives $\varphi \approx 0.1$: a plausible value, and a reminder that $q$ is
a conservative number.

\section{Method: nine stages}\label{sec:method}

\begin{figure}[H]
\centering
\begin{tikzpicture}[node distance=5mm and 6mm]
\tikzset{st/.style={text width=44mm,minimum height=13mm,align=left}}
\node[boxa,st] (m1) {\M{1} \textbf{Radar time series}\\{\scriptsize two directions, 2015--2026}};
\node[boxa,st,right=of m1] (m2) {\M{2} \textbf{Quality control}\\{\scriptsize atmosphere, GNSS tie, masks}};
\node[boxa,st,right=of m2] (m3) {\M{3} \textbf{Vertical and east--west}\\{\scriptsize combine the two directions}};
\node[boxg,st,below=of m3] (m4) {\M{4} \textbf{Time-series separation}\\{\scriptsize trend, season, change dates}};
\node[boxg,st,left=of m4] (m5) {\M{5} \textbf{Background model}\\{\scriptsize non-pumping subsidence}};
\node[boxw,st,left=of m5] (m6) {\M{6} \textbf{Sparse inversion}\\{\scriptsize few sources, cone-shaped kernel}};
\node[boxw,st,below=of m6] (m7) {\M{7} \textbf{Uncertainty}\\{\scriptsize bootstrap, stability}};
\node[box,st,right=of m7] (m8) {\M{8} \textbf{Volume to extraction}\\{\scriptsize lower bound, calibration}};
\node[box,st,right=of m8] (m9) {\M{9} \textbf{Attribution, ranking}\\{\scriptsize land use, permits, report}};
\draw[arr] (m1)--(m2); \draw[arr] (m2)--(m3); \draw[arr] (m3)--(m4); \draw[arr] (m4)--(m5);
\draw[arr] (m5)--(m6); \draw[arr] (m6)--(m7); \draw[arr] (m7)--(m8); \draw[arr] (m8)--(m9);
\end{tikzpicture}
\caption{The pipeline. Blue: measurement (\Q{1}). Green: separation (\Q{3}). Orange: attribution
(\Q{2}). White: interpretation and delivery.}\label{fig:pipeline}
\end{figure}

\subsubsection{\M{1} Radar time series}
\textbf{What.} Turn stacks of Sentinel-1 images into a displacement history for every pixel.\\
\textbf{How.} Small-baseline (SBAS) processing \cite{berardino2002}: each image is paired with its
next three or four neighbours in time, each pair gives a map of phase change, and a least-squares
solve turns the network of pairs into a time series. For the descending track we start from
ready-made LiCSAR pairs and run LiCSBAS \cite{morishita2020}, as two earlier Semarang studies did
\cite{lo2022,azeriansyah2025}. For the ascending track we order pairs from ASF's on-demand service
and run MintPy \cite{yunjun2019}. Pixels are averaged to about 100\,m. Phase $\phi$ converts to motion
along the line of sight by $d_{\mathrm{LOS}} = -\lambda\phi/4\pi$ with $\lambda = 5.55$\,cm.\\
\textbf{Output.} Two cubes (ascending, descending) of line-of-sight displacement.
\trace{serves \Q{1}; uses \D{1} \D{5}; checked by \V{2}; runs in \Sx{2}.}

\subsubsection{\M{2} Quality control and referencing}
\textbf{What.} Remove what is not ground motion, and fix the zero point.\\
\textbf{How.} (i) Correct the delay from the lower atmosphere with GACOS maps \cite{yu2018}.
(ii) Check \emph{loop closure}: for any three images, the three pairwise phases must sum to zero;
pairs that fail are unwrapping errors and are dropped. (iii) Mask pixels with low coherence
(water, fish ponds, rice fields). (iv) Reference all motion to the CSEM GNSS station, which earlier
work also used \cite{azeriansyah2025}. (v) Estimate a velocity uncertainty per pixel by resampling
the time series.\\
\textbf{Output.} Clean cubes, a validity mask, and a per-pixel uncertainty.
\trace{serves \Q{1}; uses \D{3} \D{4}; checked by \V{1} \V{2}; runs in \Sx{2}.}

\subsubsection{\M{3} Vertical and east--west motion}
\textbf{What.} A single radar direction mixes vertical and horizontal motion. Two directions untangle
them.\\
\textbf{How.} For incidence angle $\theta$ and heading $\alpha$, the line-of-sight unit vector in
(east, north, up) is
\[
\mathbf{l} = (-\sin\theta\cos\alpha,\ \ \sin\theta\sin\alpha,\ \ \cos\theta).
\]
Radar is nearly blind to north--south motion, so we solve per pixel \cite{wright2004}
\begin{equation}
\begin{pmatrix} d^{\mathrm{asc}}\\ d^{\mathrm{desc}}\end{pmatrix} =
\begin{pmatrix} l_E^{\mathrm{asc}} & l_U^{\mathrm{asc}}\\ l_E^{\mathrm{desc}} & l_U^{\mathrm{desc}}\end{pmatrix}
\begin{pmatrix}u_E\\u_U\end{pmatrix}.
\end{equation}
For $\theta \approx 34$--$39^\circ$ the vertical error is about 0.9 times the line-of-sight error and
the east--west error about 1.2 times.\\
\textbf{Output.} Vertical and east--west displacement cubes. East--west motion is a free diagnostic:
ground around a compacting bowl moves inwards.
\trace{serves \Q{1}; produces \Ox{1}; checked by \V{1}; runs in \Sx{4}.}

\subsubsection{\M{4} Time-series separation}
\textbf{What.} Split each pixel's history into long-term trend, seasonal breathing and sudden changes.\\
\textbf{How.} Three tools, from simple to rich.
\begin{itemize}
\item \emph{Robust regression} per pixel:
$u(t) = a + v\,t + c_1\sin 2\pi t + c_2\cos 2\pi t + \varepsilon$, with Huber loss so that outliers
do not bend the trend.
\item \emph{Independent component analysis} (ICA \cite{hyvarinen2000}) on the whole cube, which finds
a few shared time patterns and a map of where each one is strong. It has been used to separate
sources in radar data before \cite{ebmeier2016} and to isolate aquifer signals \cite{chaussard2014}.
\item \emph{Change-point detection} (PELT \cite{killick2012}) on the detrended series, which dates
the moments when sinking sped up or slowed down. A source that switched on or off shows up here.
\end{itemize}
\textbf{Output.} Maps of long-term rate $v$, seasonal amplitude and change dates. The seasonal part
is mostly elastic; a large trend with little seasonal recovery signals permanent compaction
\cite{hoffmann2001,smith2017}.
\trace{serves \Q{1} \Q{3}; produces \Ox{2}; uses \D{9}; checked by \V{4}; runs in \Sx{4}.}

\subsubsection{\M{5} Background model (non-pumping subsidence)}
\textbf{What.} Estimate $g(\mathbf{x})$ in Eq.~\eqref{eq:model}: the sinking that would happen without
deep pumping.\\
\textbf{How.} Two layers.
\begin{itemize}
\item A \emph{driver model}: gradient-boosted trees \cite{chen2016} that predict long-term rate from
soft-clay thickness, land age (reclaimed or not), built-up volume, distance to coast and elevation.
It is trained only on ``quiet'' pixels: those farther than $2B$ from industrial land and from any
known deep well, so that it learns background and not pumping.
\item A \emph{smooth remainder}: a low-order cosine surface that absorbs whatever long-wavelength
signal the drivers miss. It is fitted jointly with the sources in \M{6}.
\end{itemize}
Shapley values (SHAP \cite{lundberg2017}) show which driver matters where.\\
\textbf{Output.} Background map with uncertainty.
\trace{serves \Q{3}; produces \Ox{4}; uses \D{5} \D{7} \D{8}; checked by \V{4}; runs in \Sx{4}.}

\subsubsection{\M{6} Sparse source inversion}
\textbf{What.} Find the few sources $q_j$ that explain the pumping part of the velocity field.\\
\textbf{How.} Lay a grid of candidate source positions $\mathbf{c}_j$ every 500\,m. Build the matrix
$A_{ij}=k_B(\lVert\mathbf{x}_i-\mathbf{c}_j\rVert)$ and a background matrix $Z$ (driver-model
prediction plus cosine surface). Solve
\begin{equation}
\min_{\mathbf{q}\ge 0,\ \boldsymbol\beta}\;\; \frac{1}{2n}\,\big\lVert \mathbf{v}-A\mathbf{q}-Z\boldsymbol\beta\big\rVert_2^2
\;+\;\lambda\sum_j q_j .
\label{eq:lasso}
\end{equation}
This is a non-negative lasso \cite{tibshirani1996}. The constraint $q_j\ge 0$ says sources only
remove volume. The penalty $\lambda\sum q_j$ says: explain the data with as little total volume, in
as few places, as possible. Four details make it trustworthy.
\begin{enumerate}
\item \textbf{Penalty from the data's own noise.} After a first fit we take the residual map, shuffle
its Fourier phases to create 30 fake noise maps with the same spatial texture, and record the
largest source--noise correlation each one produces. $\lambda$ is set no lower than the 95th
percentile of those values. A source then enters only if it stands out from what noise alone could
do.
\item \textbf{Refit without shrinkage.} The lasso picks \emph{where}; a plain non-negative
least-squares refit on the selected positions then estimates \emph{how much}.
\item \textbf{Merging.} Neighbouring selected grid points are merged into one source at their
volume-weighted centre.
\item \textbf{Cone width $B$.} Chosen by spatial cross-validation, or better, fixed from pumping
tests or head data when available (see Section~\ref{sec:exp}).
\end{enumerate}
An optional \emph{informed mode} lowers the penalty for candidates on industrial land, as a soft
version of the well-distance penalty of Vasco et al.\ \cite{vasco2019}. The default \emph{blind mode}
uses no such prior.\\
\textbf{Output.} A source list: position and compaction volume.
\trace{serves \Q{2}; uses output of \M{3}--\M{5}; checked by \V{3} \V{5} \V{6}; runs in \Sx{4}.}

\subsubsection{\M{7} Uncertainty}
\textbf{What.} Say how sure we are of each source.\\
\textbf{How.} A \emph{spatial block bootstrap}: cut the map into 2\,km blocks, resample blocks with
replacement, and rerun the inversion. Repeating this 60--200 times gives, for each source, a 90\,\%
range for its volume and a \emph{stability} score: the share of reruns in which it reappears, in the
spirit of stability selection \cite{meinshausen2010}. Blocks are used because neighbouring pixels
share noise; resampling single pixels would give falsely narrow ranges.\\
\textbf{Output.} Volume range and stability for each source.
\trace{serves \Q{2}; produces \Ox{3}; checked by \V{7}; runs in \Sx{4}.}

\subsubsection{\M{8} From volume to extraction}
\textbf{What.} Translate compaction volume into a statement about pumping.\\
\textbf{How.} Report $q_j$ as a lower bound (Eq.~\ref{eq:phi}). Where metered extraction or head
records exist for some sources, fit $\varphi$ by regression and use it, with its spread, to convert
the other sources into estimated ranges.\\
\textbf{Output.} For each source: minimum extraction, and estimated extraction if calibration data
exist.
\trace{serves \Q{2}; uses \D{6}; checked by \V{5}; runs in \Sx{4}.}

\subsubsection{\M{9} Attribution and ranking}
\textbf{What.} Make the list usable.\\
\textbf{How.} Join each source with land use, building footprints and, if available, the permit
register. Classify it as \emph{matches a permitted well field}, \emph{no permit on record nearby}, or
\emph{unclear}. Rank by volume, then by stability. Public outputs show hot-spot cells, never company
names (Section~\ref{sec:ethics}).\\
\textbf{Output.} The audit list and report.
\trace{serves \Q{2}; produces \Ox{3} \Ox{5}; uses \D{6} \D{8}; checked by \V{6}; runs in \Sx{7} \Sx{9}.}

\section{ML/AI rigour: how we know it works}\label{sec:rigour}

\begin{inshort}
Four parts of the pipeline learn from data. Each has a baseline to beat, a split that prevents
leakage, a metric with a pass mark, and an uncertainty estimate that is itself tested.
\end{inshort}

\Needspace{27\baselineskip}
\subsection{The learning components}

\Needspace{21\baselineskip}
{\small
\begin{xltabular}{\linewidth}{@{}L{8mm}L{33mm}L{30mm}Y@{}}
\toprule
 & \textbf{Component} & \textbf{Type} & \textbf{Why it is not a black box}\\*
\midrule
\M{4} & Trend, season, change dates & Robust regression, ICA, change-point detection & Every output
  is a physical quantity (rate, amplitude, date) that can be checked against GNSS and rainfall.\\
\M{5} & Background driver model & Gradient-boosted trees with SHAP & Trained only on pixels away from
  pumping; each prediction is explained by named drivers.\\
\M{6} & Source inversion & Sparse regression (non-negative lasso) with a physics kernel & Linear and
  convex: one global optimum, every source is a visible term in Eq.~\eqref{eq:model}.\\
\M{7} & Uncertainty & Block bootstrap, stability selection & Ranges are tested for coverage on
  synthetic truth (\V{7}).\\
\bottomrule
\end{xltabular}}

\subsection{Splits and leakage rules}

Map data break the usual assumption that samples are independent: a pixel's neighbours are almost
copies of it. A random train/test split would then report excellent scores for a model that has
learned nothing transferable \cite{roberts2017}. We therefore fix these rules.

\begin{itemize}
\item \textbf{Spatial block cross-validation.} Pixels are grouped into 2\,km blocks; whole blocks are
held out. Used for choosing $B$, $\lambda$ and the driver model's settings.
\item \textbf{Temporal hold-out.} All fitting uses 2015--2023. The years 2024--2026 are kept back.
A good model must predict the later cumulative motion better than simply extending each pixel's own
trend.
\item \textbf{No coordinates as features.} The driver model never sees latitude or longitude.
Otherwise it would memorise where the bowls are.
\item \textbf{No target-derived features.} Nothing computed from the velocity map (for example
``distance to nearest bowl'') may enter \M{5}.
\item \textbf{Quiet-pixel training.} \M{5} is trained only where pumping influence is unlikely, and
evaluated on held-out quiet blocks.
\item \textbf{Frozen test set.} Calibration and validation wells (\D{6}) are split once, by well, and
the validation half is not looked at until the pipeline is final.
\end{itemize}

\Needspace{14\baselineskip}
\Needspace{14\baselineskip}
\subsection{Validation checks and pass marks}

\Needspace{8\baselineskip}
{\small
\begin{xltabular}{\linewidth}{@{}L{8mm}L{37mm}Y L{37mm}@{}}
\caption{Validation checks. Pass marks are ours, set from what earlier studies achieved and from the
synthetic experiment.}\label{tab:valid}\\*
\toprule
 & \textbf{Check} & \textbf{How it is measured} & \textbf{Pass mark}\\*
\midrule
\V{1} & Radar against GNSS & Vertical rate at InaCORS stations, root-mean-square difference & $\le 8$\,\mmyr\
  (earlier work: 3--9 \cite{azeriansyah2025,janur2025})\\
\V{2} & Internal consistency & Loop-closure failures; agreement of ascending and descending where motion is
  vertical & $<5$\,\% of pairs dropped\\
\V{3} & Synthetic recovery & Known sources planted in a simulated map (Section~\ref{sec:exp}) & Recall and
  precision $\ge 80$\,\%; median volume error $\le 25$\,\%\\
\V{4} & Out-of-sample fit & Block cross-validation error; temporal hold-out against trend extension & Beats
  both baselines\\
\V{5} & Wells: heads & Rank correlation between inferred compaction rate and head decline at monitoring
  wells & $\ge 0.6$\\
\V{6} & Wells: locations & Share of registered high-capacity well clusters that lie within 1\,km of a
  detected source & $\ge 70$\,\%\\
\V{7} & Honest ranges & Share of true volumes inside the stated 90\,\% range (synthetic) & 80--95\,\%\\
\bottomrule
\end{xltabular}}

\subsection{Baselines and ablations}

A method is only as convincing as what it is compared with. Three baselines are run on every data
set.

\begin{itemize}
\item \textbf{Bowl integration.} Threshold the subsidence map, treat each connected bowl as one source,
integrate its volume. This is what a careful analyst would do by hand.
\item \textbf{Smooth inversion.} The same kernel and data as ours, but with a smoothness (L2) penalty
instead of the sparsity penalty, in the spirit of \cite{azeriansyah2025}.
\item \textbf{Trend extension.} For the temporal hold-out: predict 2024--2026 by extending each
pixel's 2015--2023 trend.
\end{itemize}

Ablations switch off one ingredient at a time: one radar direction instead of two; no background
model; no noise-calibrated penalty; no refit; Gaussian instead of Bessel kernel; blind against
informed mode. Each ablation reports the metrics of Table~\ref{tab:valid}.

\subsection{Reproducibility}

Fixed random seeds; all settings in one versioned configuration file; data versions tracked by
content hash; every run logged with its parameters, metrics and figures; the full pipeline runs from
one command inside a container. Section~\ref{sec:mlops} gives the tooling. The synthetic experiment
in the next section follows these rules and its code ships with this document.

\subsection{Stretch goal: a learned inversion}

The forward model \eqref{eq:model} is cheap to simulate. That makes it possible to train a
convolutional network (U-Net \cite{ronneberger2015}) on tens of thousands of simulated maps to
predict source maps directly, as has been done for detecting volcanic deformation in radar data
\cite{anantrasirichai2018}. It would be benchmarked against \M{6} on the same synthetic test set.
It is not on the critical path: the convex inversion is simpler to defend in front of a jury, and
its failure modes are easier to explain.
